\section{Giới thiệu}
\label{sec:introduction}

Trong những năm gần đây, các large language model (LLM) thương mại đã tiến từ chỗ hỗ
trợ soạn thảo từng phần sang đảm đương gần trọn quy trình instructional systems
design (ISD): các hệ thống dựa trên GPT-4 và những model cùng lớp có thể sinh được
lesson plan, bộ câu hỏi đánh giá, thậm chí trọn gói học liệu của một môn học từ một
yêu cầu thiết kế duy
nhất~\cite{yao2025instructionalagents,zhang2025eduplanner,agentlesson2026,zeng2026coma,hu2024teachingplan}.
Chất lượng ở cấp thành phần đã được kiểm chứng qua nhiều nghiên cứu so sánh:
learning objective tương đương bản do giảng viên soạn~\cite{doyle2025objectives},
câu hỏi trắc nghiệm ngang bản người viết~\cite{doughty2024mcq}, giáo án không phân
biệt được với bản của đồng nghiệp khi chấm
blind~\cite{cooper2026indistinguishable}, và hiệu quả lớp học không thua
kém~\cite{karaman2024chatgptlesson}; với vòng lặp mô phỏng--tinh chỉnh, chất lượng
đạt mức của giáo viên giàu kinh nghiệm~\cite{hu2025teachingsim}. Các công cụ này
cũng đã đi vào thực hành của instructional
designer~\cite{yang2025idpractice,lessonplanner2024}.

Tuy nhiên, việc triển khai các hệ thống này tại cơ sở giáo dục gặp trở ngại vì
chúng phụ thuộc vào API thương mại. Một quy trình ISD hoàn chỉnh cần nhiều lời gọi model nên
chi phí tích luỹ lớn; khảo sát trong giáo dục đại học ghi nhận ràng buộc tài chính
đang làm chậm việc ứng dụng AI~\cite{robert2025educause}, còn dữ liệu người học
phải đi qua hạ tầng bên thứ ba trong khi vấn đề riêng tư được đánh giá là chưa xử
lý thoả đáng~\cite{yan2024ethical}. Nhiều đại học đã dựng gateway riêng để kiểm
soát ngân sách và dữ liệu~\cite{mithun2025aiverde}; tự vận hành model open-weight
cũng rẻ hơn API khi khối lượng sử dụng đủ lớn~\cite{pan2025onpremise}. Như vậy,
chi phí và an toàn dữ liệu được ghi nhận là rào cản triển khai tại các cơ sở giáo
dục, và hướng small language model (SLM) open-weight chạy on-premise đang được
quan tâm: hạ tầng LLM mở cho đại học đào tạo từ xa~\cite{zesch2024flexi}, đánh giá
model tự host cho trường học~\cite{chan2026bigger}, khung biên soạn nội dung chạy
cục bộ vì riêng tư và chi phí~\cite{reza2025smallmodels}.

Khi dùng SLM, model nền yếu hơn, nên chất lượng đầu ra phụ thuộc nhiều hơn vào
harness bao quanh, tức là cách công việc được phân rã, điều phối và kiểm soát chất
lượng. Model nhỏ yếu khi dùng tool nhưng khung đa tác tử có phân rã thu hẹp khoảng
cách với model lớn~\cite{shen2024multillmagent}; scaffolding dạng agent giúp model
0.5B--3B đạt gần chất lượng của nhóm lớn hơn~\cite{kang2025agentdistill}; SLM cũng
được lập luận là đủ cho phần lớn agentic
workload~\cite{belcak2025slm,sharma2025slmagents}. Tuy nhiên, độ nhạy với harness
không đơn điệu giữa các mức kích thước model~\cite{cho2026harness}, nên quan hệ
giữa chất lượng, harness và kích thước model cần được đo lường thay vì giả định.

Từ bối cảnh trên, khoảng trống nghiên cứu trung tâm hiện ra: các hệ tác tử ISD và
lesson-plan hiện có đều được xây dựng và đánh giá trên model thương mại mạnh, còn
các công trình gần nhất mỗi công trình chỉ chạm một phần của bài
toán~\cite{jaldi2026slmassessment,zheng2025lessonplanlm,reza2025smallmodels,lee2026educlaw,wei2026slmeducation}.
Theo khảo sát của chúng tôi, chưa có công trình nào cho thấy một hệ tác tử duy
trì được chất lượng ISD end-to-end trên model nhỏ mà không cần fine-tune, cũng
như chưa có đánh giá hệ thống nào về chất lượng đó trên một dải kích thước model.

Đồng thời, khoảng trống đó cũng làm nảy sinh một vấn đề khác liên quan đến phương pháp đo lường chất lượng,
bởi hai công cụ thông dụng đều có giới hạn. Rubric chấm bằng LLM-judge --- cách của các
benchmark ISD và giáo dục hiện
nay~\cite{isdagentbench,edubench2025,tutorbench2025,mathtutorbench2025} --- mang
các bias đã được ghi nhận có hệ
thống~\cite{zheng2023judging,panickssery2024llm,norman2026reliability}, và ở đây
họ model chấm còn trùng với họ model sinh. Đánh giá bằng chuyên gia con người thì
đáng tin hơn nhưng không scale cho một ladder gồm nhiều kích thước, nhiều run và
nhiều hệ. Quan trọng hơn cả, constructive alignment đặt chất lượng của một thiết
kế ở quan hệ giữa mục tiêu, hoạt động và đánh giá~\cite{biggs1996enhancing}: một
thiết kế có thể đọc trôi chảy, nhìn qua có vẻ tốt, nhưng vẫn sai lệch nghiêm trọng về
alignment. Không cách đo nào ở trên làm tường minh được các tín hiệu alignment
ấy. Các phương pháp đo alignment tự động hiện có cũng chỉ đo sự ăn khớp
với một bộ chuẩn hoặc tham chiếu soạn
sẵn~\cite{porter2002measuring,webb1997criteria,xu2025itemalignment,silva2026lessonbench},
chưa đo alignment nội tại của một thiết kế hoàn toàn do máy sinh.

Mục tiêu của nghiên cứu này là duy trì chất lượng thiết kế ISD khi thu nhỏ model
bằng một harness đa tác tử có trạng thái alignment tường minh
(\Cref{sec:method}), đồng thời làm tường minh các tín hiệu alignment của thiết kế
được sinh bằng một giao thức đo không dùng LLM (\Cref{sec:protocol}). Mục tiêu đó
được cụ thể hoá thành hai câu hỏi nghiên cứu:
\begin{itemize}
  \item \textbf{RQ1 (duy trì chất lượng theo thang LLM-as-judge).} Khi model nền
  nhỏ dần trên ladder Qwen 0.8B--9B, AlignmentGraph-ISD có duy trì được chất lượng
  thiết kế đo bằng LLM-as-judge của ISD-Agent-Bench so với sáu benchmark agent
  không: hệ đứng ở đâu tại từng kích thước, và khoảng cách có nới rộng khi model
  nhỏ đi không?
  \item \textbf{RQ2 (duy trì chất lượng theo thang tín hiệu alignment tường minh).} Khi được làm tường
  minh bằng giao thức đo không dùng LLM, các tín hiệu alignment nội tại của thiết
  kế do máy sinh thay đổi thế nào theo kích thước model, và khoảng cách so sánh có
  nới rộng khi model nhỏ đi không?
\end{itemize}

Bài báo có hai đóng góp chính:
\begin{enumerate}
  \item \textbf{AlignmentGraph-ISD}: một harness đa tác tử tổ chức việc sinh
  chương trình ADDIE quanh một alignment graph, tức trạng thái alignment tường
  minh, có thể kiểm tra và audit. Cuối mỗi lần chạy, toàn bộ graph được xuất ra để
  instructional designer xem xét từng liên kết mục tiêu--hoạt động--đánh giá và
  can thiệp khi cần, theo nguyên tắc teacher-in-the-loop.
  \item \textbf{Giao thức đánh giá constructive alignment tự động}, mở rộng
  ISD-Agent-Bench: không dùng LLM trong đo lường, mỗi điểm số là một hàm cố định
  của văn bản, kết quả là panel bảy tín hiệu ngang hàng thuộc hai họ công cụ đo
  độc lập về cơ chế. Giao thức operationalize các tiêu chí đã được thiết lập thay
  vì đề xuất construct
  mới~\cite{biggs1996enhancing,porter2002measuring,webb1997criteria}; phạm vi
  diễn giải giới hạn ở các tín hiệu đo được.
\end{enumerate}

Phần còn lại của bài báo được tổ chức như sau: \Cref{sec:related-work} trình bày
các công trình liên quan; \Cref{sec:method} mô tả AlignmentGraph-ISD và giao thức
đánh giá (\Cref{sec:protocol}); \Cref{sec:experimental-methodology} mô tả thiết kế
thực nghiệm; \Cref{sec:results} báo cáo kết quả; \Cref{sec:discussion} thảo luận
kết quả và các nguy cơ đối với tính hợp lệ; \Cref{sec:conclusion} kết luận bài
báo.
